Most unsupervised methods for representation learning can be categorized as either generative or discriminative~\cite{doersch2015unsupervised, SimCLR}. 
Generative approaches to representation learning build a distribution over data and latent embedding and use the learned embeddings as image representations. Many of these approaches rely either on auto-encoding of images~\cite{vincent2008extracting, KingmaVAE, RezendeVAE} or on adversarial learning~\cite{GAN}, jointly modelling data and representation~\cite{donahue2016adversarial, DumoulinALI, BigBiGAN, BrockBigGAN}. Generative methods typically operate directly in pixel space. This however is computationally expensive, and the high level of detail required for image generation may not be necessary for representation learning.

Among discriminative methods, contrastive methods~\cite{MOCO, CPC, CPCv2, DIM, AMDIM, CMC, PIRL, li2020prototypical}
currently achieve state-of-the-art performance in self-supervised learning~\cite{MOCOv2, SimCLR,chen2020big,Tian2020WhatMF}. Contrastive approaches avoid a costly generation step in pixel space by bringing representation of different views of the same image closer (`positive pairs'), and spreading representations of views from different images (`negative pairs') apart~\cite{wu2018unsupervised, doersch2017multi}. Contrastive methods often require comparing each example with many other examples to work well~\cite{MOCO, SimCLR} prompting the question of whether using negative pairs is necessary. 

\texttt{DeepCluster}~\cite{caron2018DeepCF} partially answers this question. It uses bootstrapping on previous versions of its representation to produce targets for the next representation; it clusters data points using the prior representation, and uses the cluster index of each sample as a classification target for the new representation. While avoiding the use of negative pairs, this requires a costly clustering phase and specific precautions to avoid collapsing to trivial solutions. 

Some self-supervised methods are not contrastive but rely on using auxiliary handcrafted prediction tasks to learn their representation. In particular, relative patch prediction \cite{doersch2015unsupervised,doersch2017multi}, colorizing gray-scale images~\cite{Zhang2016ColorfulIC, larsson2016learning}, image inpainting~\cite{pathak2016context}, image jigsaw puzzle~\cite{noroozi2016unsupervised}, image super-resolution~\cite{ledig2017photo}, and geometric transformations~\cite{dosovitskiy2014discriminative, gidaris2018unsupervised} have been shown to be useful. Yet, even with suitable architectures~\cite{kolesnikov2019revisiting}, these methods are being outperformed by contrastive methods~\cite{MOCOv2, SimCLR, Tian2020WhatMF}.

Our approach has some similarities with \textit{Predictions of Bootstrapped Latents} (\texttt{PBL}, \cite{guo2020bootstrap}), a self-supervised representation learning technique for reinforcement learning (RL). \texttt{PBL} jointly trains the agent's history representation and an encoding of future observations. The observation encoding is used as a target to train the agent's representation, and the agent's representation as a target to train the observation encoding. Unlike \texttt{PBL}, \algo uses a slow-moving average of its representation to provide its targets, and does not require a second network.


The idea of using a slow-moving average target network to produce stable targets for the online network was inspired by  deep RL \cite{DQNNature, mnih2016asynchronous, hessel2018rainbow, van2018deep}. Target networks stabilize the bootstrapping updates provided by the Bellman equation, making them appealing to stabilize the bootstrap mechanism in \algo. While most RL methods use fixed target networks, \algo uses a weighted moving average of previous networks (as in \cite{DDPG}) in order to provide smoother changes in the target representation. % Related ideas have been explored in the deep learning literature under the name \textit{distillation}\cite{hinton2015distilling,Zhang_2019_ICCV}.


In the semi-supervised setting~\cite{chapelle2009semi,zhu2009introduction}, an unsupervised loss is combined with a classification loss over a handful of labels to ground the training~\cite{laine2016temporal,tarvainen2017mean,kingma2014semi,rasmus2015semi,berthelot2019mixmatch,miyato2018virtual,Berthelot2020ReMixMatchSL,sohn2020fixmatch}. Among these methods, \textit{mean teacher} (\MT) \cite{tarvainen2017mean} also uses a slow-moving average network, called \textit{teacher}, to produce targets for an online network, called \textit{student}. % and predicts them via an $\ell_2$-loss.
% A consistency loss, which is an $\ell_2$-loss between the softmax predictions of the teacher and the student, is added to the classification loss.
An $\ell_2$ consistency loss between the softmax predictions of the teacher and the student is added to the classification loss.
While \cite{tarvainen2017mean} demonstrates the effectiveness of \MT in the semi-supervised learning case, in~\Cref{sec:mean_teacher} we show that a similar approach collapses when removing the classification loss. In contrast, \algo introduces an additional predictor on top of the online network, which prevents collapse. 

Finally, in self-supervised learning, \moco~\cite{MOCO} uses a slow-moving average network (\emph{momentum encoder}) to maintain consistent representations of negative pairs drawn from a memory bank. Instead, \algo uses a moving average network to produce prediction targets as a means of stabilizing the bootstrap step. We show in \Cref{subsec:simclr_to_pbl} that this mere stabilizing effect can also improve existing contrastive methods.

